\pdfvariable suppressoptionalinfo 611
\providecommand{\CRevisionRound}{2}
\documentclass[11pt]{article}
\usepackage[a4paper,margin=22mm]{geometry}
\usepackage{amsmath,amssymb,amsthm,booktabs,microtype}
\usepackage[T1]{fontenc}
\usepackage{lmodern}
\usepackage[hidelinks,hypertexnames=false]{hyperref}
\newtheorem{theorem}{Theorem}
\newcommand{\dd}{\,\mathrm d}
\title{Rayleigh Spherical-Cavity Collapse: Exact Beta Clock and the Terminal \(2/5\) Singularity}
\author{HCS Research Program}
\date{28 August 2026\quad (revision \CRevisionRound)}
\begin{document}
\maketitle
\begin{abstract}
We give a source-local all-parameter atlas for the inviscid spherical-cavity
Rayleigh equation with constant pressure difference.  A weighted first
integral closes the positive, zero, and negative pressure branches.  On the
physical collapse branch the lifetime is an exact Beta constant, while the
terminal radius, wall speed, wall acceleration, volume, finite liquid kinetic
energy, and \(L^p\) thresholds follow from one endpoint expansion.  The
zero-radius face is kept outside the classical phase space.  This is a
hydrodynamic theorem and reproducibility certificate; it is not an arithmetic
determinant or a Hilbert--Pólya construction.
\end{abstract}

\section{Model and exact first integral}
Let \(R(t)>0\) be the radius of an empty spherical cavity in an ideal
incompressible liquid.  We freeze
\begin{equation}
 R\ddot R+\frac32\dot R^2=-\frac{\Pi}{\rho},
 \qquad R(0)=R_0>0,\quad \dot R(0)=0,
 \tag{1}
\end{equation}
where \(\rho>0\) and \(\Pi=P_\infty-P_v\) is constant.  The physical
collapse sign is \(\Pi>0\); the other signs are retained as mathematical
controls.  Multiplication by \(2R^2\dot R\) gives
\begin{equation}
 R^3\dot R^2+\frac{2\Pi}{3\rho}R^3
 =\frac{2\Pi}{3\rho}R_0^3. \tag{2}
\end{equation}

\begin{theorem}[complete sign and endpoint atlas]\label{thm:atlas}
For \(\Pi>0\), put \(a=\sqrt{2\Pi/(3\rho)}\).  The positive-radius forward
solution is strictly decreasing until its finite collapse time
\begin{equation}
 T_c=\frac{R_0}{a}\int_0^1\frac{x^{3/2}}{\sqrt{1-x^3}}\dd x
 =\frac{R_0}{3a}B\!\left(\frac56,\frac12\right)
 =0.914681356501962\ldots R_0\sqrt{\frac{\rho}{\Pi}}. \tag{3}
\end{equation}
Writing \(\delta=T_c-t\),
\begin{equation}
 R=C\delta^{2/5}(1+O(\delta^{6/5})),\quad
 \dot R=-\frac25C\delta^{-3/5}(1+O(\delta^{6/5})),\quad
 \ddot R=-\frac6{25}C\delta^{-8/5}(1+O(\delta^{6/5})),
 \tag{4}
\end{equation}
where \(C=R_0^{3/5}(5a/2)^{2/5}\).  For \(\Pi=0\),
\(R(t)=R_0\).  For \(\Pi<0\), the solution is strictly increasing, has no
finite singular time, and \(R(t)\sim\sqrt{2|\Pi|/(3\rho)}\,t\).
Moreover, with \(V=(4\pi/3)R^3\),
\begin{equation}
 K=2\pi\rho R^3\dot R^2=\frac{4\pi}{3}\Pi(R_0^3-R^3),\quad
 U=\frac{4\pi}{3}\Pi R^3,\quad E=K+U=\frac{4\pi}{3}\Pi R_0^3, \tag{5}
\end{equation}
and near collapse \(V\sim(4\pi/3)C^3\delta^{6/5}\).  Finally,
\[
 \dot R\in L^p(T_c-\varepsilon,T_c)\ \Longleftrightarrow\ p<\frac53,
 \qquad
 \ddot R\in L^p(T_c-\varepsilon,T_c)\ \Longleftrightarrow\ p<\frac58.
 \tag{6}
\]
\end{theorem}

\begin{proof}
Equation (2) follows directly from (1).  On the collapse branch its negative
square-root sign gives
\(\dot R=-a\sqrt{(R_0/R)^3-1}\), and separation gives (3).  The substitution
\(y=x^3\) produces \(\frac13B(5/6,1/2)\).  At the endpoint,
\[
 \int_0^x\frac{u^{3/2}}{\sqrt{1-u^3}}\dd u
 =\frac25x^{5/2}+\frac1{11}x^{11/2}+O(x^{17/2}),
\]
which inverts to (4).  The same powers prove (6).  The positive and negative
signs in (2) give the monotonicity and the expansion asymptotic.  Substitution
of (2) into the kinetic term proves (5) and the volume statement.
\end{proof}

\section{Lagrangian and geometric boundary}
The equation is Euler--Lagrange on \(R>0\) for
\begin{equation}
 L_{\rm phys}=2\pi\rho R^3\dot R^2-\frac{4\pi}{3}\Pi R^3. \tag{7}
\end{equation}
Indeed its residual is
\(4\pi\rho R^2[R\ddot R+(3/2)\dot R^2+\Pi/\rho]\).  Thus the wall speed
diverges while the liquid kinetic energy tends to the finite value
\(4\pi\Pi R_0^3/3\), and the volume flux tends to zero.  The dimensionless
 profile is an inverse incomplete-Beta function; this notation is exact but is
not an elementary closed form.  It is source-local explicit solvability only:
the source Beta clock is not target continuation/divisor/counting law and is
not an A3 analytic-structure match.  The face \(R_0=0\) is not a
positive-radius initial state, so no continuation through \(R=0\) is claimed.

\ifnum\CRevisionRound>0
\section{Boundary controls and reproducible ledger}
The certificate samples five collapse rows, two equilibria, three expansion
rows, and three zero-radius controls.  The producer uses endpoint-stable
quadratures; an independent checker uses hypergeometric primitives.  SymPy
verifies the first integral, Beta substitution, Euler--Lagrange residual,
Puiseux coefficients, volume/energy identities, and both integrability
thresholds.  A byte replay and repaired/stale-hash mutation suite are part of
the release contract.
\fi

\ifnum\CRevisionRound>1
\section{Route-A scope and audit}
The strict tuple is
\begin{center}\ttfamily
(A0\_FAIL, A1\_FAIL, A2\_FAIL, A3\_FAIL, A4\_FORMAL\_HINT).
\end{center}
Thus \texttt{overall=ROUTE\_A\_REJECTED} and
\texttt{route\_b\_invocation\_allowed=false}.
The monotone collapse trajectory has no primitive periodic-orbit owner,
arithmetic clock, target determinant, or same-clock self-adjoint lift.  The
incomplete-Beta branch and terminal Puiseux law are source-local explicit
solvability only; the source Beta clock is not target continuation/divisor/counting
law and is not an A3 target match.  The registered scope is
\begin{center}\ttfamily NO\_BAD\_EULER\_OR\_ROOT\_NUMBER\end{center}
No prime table, zero table, local factor, root number, automorphy statement,
functional equation, or Hilbert--Pólya operator is introduced.
\fi

\section*{Source note}
The ideal empty-cavity equation originates with Rayleigh~\cite{rayleigh1917}.
The bubble-dynamics context is reviewed by Plesset and Prosperetti
\cite{plesset1977}; the terminal \(2/5\) singularity and physical correction
mechanisms are discussed by Brenner, Hilgenfeldt, and Lohse~\cite{brenner2002}.
The dimensionless profile and collapse approximation are recorded by
Obreschkow, Bruderer, and Farhat~\cite{obreschkow2012}.  These references are
source metadata, not a priority or arithmetic claim.

\begin{thebibliography}{9}
\bibitem{rayleigh1917} Lord Rayleigh, ``VIII. On the pressure developed in a
liquid during the collapse of a spherical cavity,'' \emph{The London,
Edinburgh, and Dublin Philosophical Magazine and Journal of Science} 34(200),
94--98 (1917). DOI: \href{https://doi.org/10.1080/14786440808635681}{10.1080/14786440808635681}.
\bibitem{plesset1977} M. S. Plesset and A. Prosperetti, ``Bubble Dynamics and
Cavitation,'' \emph{Annual Review of Fluid Mechanics} 9, 145--185 (1977).
DOI: \href{https://doi.org/10.1146/annurev.fl.09.010177.001045}{10.1146/annurev.fl.09.010177.001045}.
\bibitem{brenner2002} M. P. Brenner, S. Hilgenfeldt, and D. Lohse,
``Single-bubble sonoluminescence,'' \emph{Reviews of Modern Physics} 74,
425--484 (2002). DOI: \href{https://doi.org/10.1103/RevModPhys.74.425}{10.1103/RevModPhys.74.425}.
\bibitem{obreschkow2012} D. Obreschkow, M. Bruderer, and M. Farhat,
``Analytical approximations for the collapse of an empty spherical bubble,''
\emph{Physical Review E} 85, 066303 (2012). DOI:
\href{https://doi.org/10.1103/PhysRevE.85.066303}{10.1103/PhysRevE.85.066303}.
\end{thebibliography}

\section*{Declarations}
\textbf{Scope.} \texttt{NO\_BAD\_EULER\_OR\_ROOT\_NUMBER}; no target arithmetic
or operator claim.  \textbf{Data and code.} The theorem, ledger, independent
checks, and build instructions are released with HCS-C219.
\textbf{AI-use disclosure.} Generative tools assisted drafting and code
generation; the internal artifact chain checked the displayed claims and
metadata.  This is not external peer review.
\end{document}
